\section{Prediction-Powered Fusion Learning for CATE}
\label{sec:cate}

We develop two learners for the RCT-population CATE \(\tau(x)\). The primary
\emph{DRF-Learner} regresses an RCT-valid AIPW pseudo-outcome. The secondary
\emph{RF-Learner} uses a normalized residual loss with matching curvature.
All learned objects are fixed before fresh fitting or validation observations
are used, as required by Assumption~\ref{ass:honest-separation}.

\subsection{Target-preserving prediction-powered losses}
\label{sec:cate-loss}

For fixed \(t\), define
\begin{align}
\widehat L_{\mathrm{DRF}}(t;\lambda,\omega,r)
&\triangleq \mathbb P_R^{\mathrm{fit}}\{Z_\lambda-t(X)\}^2
+\omega\left[
\mathbb P_O^{\mathrm{fit}}\{r(X)(g(X)-t(X))^2\}
-\mathbb P_R^{\mathrm{fit}}\{(g(X)-t(X))^2\}
\right].
\label{eq:dr-fusion-loss}
\end{align}
Assume \(Z_\lambda,t\in L^2(P_R)\), and let
\(G_t(X)\triangleq\{g(X)-t(X)\}^2\) satisfy
\(G_t\in L^1(P_R)\) and \(r_0G_t\in L^1(P_O)\).

\begin{proposition}[Exact DRF target]
\label{prop:cate-target}
Under Assumptions~\ref{ass:trial} and \ref{ass:honest-separation}, suppose
either Assumption~\ref{ass:common-marginal} holds and \(r_0=1\), or
Assumption~\ref{ass:exact-transport} holds and
\(r_0=dP_R^X/dP_O^X\). Use \(r=r_0\). Then
\begin{align}
\mathbb E\{\widehat L_{\mathrm{DRF}}(t;\lambda,\omega,r_0)\}
&=Q_{\mathrm{DRF}}(\lambda)+\mathbb E_R\{t(X)-\tau(X)\}^2,
\nonumber\\
Q_{\mathrm{DRF}}(\lambda)
&\triangleq\mathbb E_R\{Z_\lambda-\tau(X)\}^2.
\label{eq:cate-target}
\end{align}
Thus the population minimizers over any fixed class \(\mathcal T\) are the
\(L_2(P_R^X)\) projections of \(\tau\). Under the common covariate marginal,
\(r_0=1\). The expected loss is flat in \(\omega\).
\end{proposition}

The target identity does not make the numerical loss orthogonal to every
nuisance. The relevant result concerns the projection score.

\begin{proposition}[Partial-path DRF score orthogonality]
\label{prop:drf-score-orthogonality}
For fixed \(v,g,r_0\), define
\begin{align*}
\Psi_{\mathrm{DRF}}(t,q)[v]
&\triangleq\mathbb E_R[v(X)\{\varphi(V;q)-t(X)\}].
\end{align*}
With known \(e_0\), its pathwise derivative under a perturbation of \(q\) is
zero. The derivative of the numerical squared loss need not be zero. Moreover,
\begin{align}
D_rL_{\mathrm{DRF}}[h]
&=\omega\mathbb E_O[h(X)\{g(X)-t(X)\}^2],
\label{eq:drf-ratio-derivative}
\end{align}
which is generally nonzero.
\end{proposition}

This is a partial-path statement holding \((g,r_0)\) fixed, not joint
orthogonality when shared OBS regressions generate both \(q_O\) and \(g\).
Known RCT propensity is load-bearing for arbitrary fused \(q_\lambda\).

For a fixed marginal-outcome function \(m\), define
\begin{align*}
\kappa(A,X)&\triangleq
\frac{\{A-e_0(X)\}^2}{e_0(X)\{1-e_0(X)\}},
&\mathbb E_R(\kappa\mid X)&=1,\\
\ell_R^{\mathrm{norm}}(t;m)&\triangleq
\frac{\{Y-m(X)-\{A-e_0(X)\}t(X)\}^2}
{e_0(X)\{1-e_0(X)\}}.
\end{align*}
The normalization is necessary because the raw R-loss induces an
\(e_0(X)\{1-e_0(X)\}\)-weighted projection in a restricted class. Define
\begin{align}
\widehat L_{\mathrm{RF}}(t;m,\omega,r)
&\triangleq\mathbb P_R^{\mathrm{fit}}\ell_R^{\mathrm{norm}}(t;m)
+\omega\left[
\mathbb P_O^{\mathrm{fit}}\{r(X)G_t(X)\}
-\mathbb P_R^{\mathrm{fit}}\{\kappa(A,X)G_t(X)\}
\right].
\label{eq:rf-fusion-loss}
\end{align}
Assume the normalized residual square is integrable and \(\tau\in L^2(P_R)\),
and define
\begin{align*}
C_R(m)&\triangleq\mathbb E_R\left[
\frac{\{Y-m(X)\}^2}{e_0(X)\{1-e_0(X)\}}-\tau(X)^2\right].
\end{align*}

\begin{theorem}[Exact normalized RF target]
\label{thm:rf-target}
Under the assumptions of Proposition~\ref{prop:cate-target},
\begin{align}
\mathbb E\{\widehat L_{\mathrm{RF}}(t;m,\omega,r_0)\}
&=C_R(m)+\mathbb E_R\{t(X)-\tau(X)\}^2.
\label{eq:rf-target}
\end{align}
Hence the population minimizers over a fixed class are the same unweighted
\(L_2(P_R^X)\) projections of \(\tau\).
\end{theorem}

\begin{proposition}[Partial-path RF score orthogonality]
\label{prop:rf-score-orthogonality}
For fixed \(v,g,r_0\), define
\begin{align*}
\Psi_{\mathrm{RF}}(t,m)[v]&\triangleq\mathbb E_R\left[
v(X)\frac{A-e_0(X)}{e_0(X)\{1-e_0(X)\}}
\{Y-m(X)-\{A-e_0(X)\}t(X)\}\right].
\end{align*}
Its pathwise derivative with respect to \(m\) is zero. For an estimated
propensity, propensity orthogonality holds only locally at
\((t,m,e)=(\tau,m_0,e_0)\), where
\(m_0(X)\triangleq\mathbb E_R(Y\mid X)\). Ratio perturbations remain
nonorthogonal as in Equation~\eqref{eq:drf-ratio-derivative}.
\end{proposition}

Define
\begin{align*}
m_R(x)&\triangleq e_0(x)q_{R,1}(x)+\{1-e_0(x)\}q_{R,0}(x),\\
m_O(x)&\triangleq e_0(x)q_{O,1}(x)+\{1-e_0(x)\}q_{O,0}(x),\\
m_\lambda(x)&\triangleq(1-\lambda)m_R(x)+\lambda m_O(x).
\end{align*}
Theorem~\ref{thm:rf-target} holds for every fixed \(m_\lambda\), so
\(\lambda\) is target-preserving but may change residual noise.

\subsection{Final learner-risk coefficients and honest selection}
\label{sec:cate-risk-selection}

For \(j\in\{\mathrm{DRF},\mathrm{RF}\}\), let a prespecified finite class
\(\mathfrak C_j\) index \((\lambda,\omega,p,\rho)\) and all other fitting
choices. For candidate \(c\), define
\begin{align}
\mathcal R_{\mathrm{fit}}^j(c)&\triangleq
\mathbb E_R[\{\widehat t_c^j(X)-\tau(X)\}^2],
\label{eq:sample-cate-risk}\\
\mathfrak R_j(c)&\triangleq
\mathbb E_{\mathrm{fit}}\{\mathcal R_{\mathrm{fit}}^j(c)\}.
\label{eq:algorithm-cate-risk}
\end{align}
When the minimum is attained, define the final learner-risk candidate and its
two fusion coordinates by
\begin{align}
c_j^\star&\in\operatorname*{arg\,min}_{c\in\mathfrak C_j}\mathfrak R_j(c),
&
(\lambda_j^\star,\omega_j^\star)&\triangleq
(\lambda_{c_j^\star},\omega_{c_j^\star}).
\label{eq:final-risk-oracle}
\end{align}
There is no universal closed form for this algorithm-induced oracle. ATE
variance, pseudo-outcome-noise, and fixed-loss-variance oracles are different
objects, as detailed in the Appendix.

Let \(\mathfrak C_{j,R}\subseteq\mathfrak C_j\) contain the identical
RCT-only learners, with OBS channels set to zero.
\begin{proposition}[Candidate-class oracle dominance]
\label{prop:cate-candidate-dominance}
For every realized fitting sample,
\begin{align}
\inf_{c\in\mathfrak C_j}\mathcal R_{\mathrm{fit}}^j(c)
&\leq\inf_{c\in\mathfrak C_{j,R}}\mathcal R_{\mathrm{fit}}^j(c).
\label{eq:sample-oracle-dominance}
\end{align}
The analogous expected-training inequality also holds. Strict improvement
requires a fusion candidate with strictly smaller risk. This oracle comparison
does not imply finite-sample dominance for estimated coefficients.
\end{proposition}

Use one shared RCT-valid \(Z_{\mathrm{ref}}\), independent of validation. For
\(G_c(X)\triangleq\{g(X)-\widehat t_c^j(X)\}^2\), define
\begin{align}
\widehat{\mathcal R}^{\mathrm{ref}}_c
&\triangleq\mathbb P_{R,\mathrm{val}}
\{Z_{\mathrm{ref}}-\widehat t_c^j(X)\}^2
+\omega_c\left[
\mathbb P_{O,\mathrm{val}}\{r^\dagger G_c\}
-\mathbb P_{R,\mathrm{val}}G_c\right].
\label{eq:common-reference-validation}
\end{align}
Under Assumptions~\ref{ass:trial} and \ref{ass:honest-separation}, use
\(r^\dagger=r_0=1\) under Assumption~\ref{ass:common-marginal}, or
\(r^\dagger=r_0\) under Assumption~\ref{ass:exact-transport}. Then
\begin{align}
\mathbb E_{\mathrm{val}}\widehat{\mathcal R}^{\mathrm{ref}}_c
&=Q_{\mathrm{ref}}+\mathcal R_{\mathrm{fit}}^j(c),
\label{eq:common-reference-identity}
\end{align}
where \(Q_{\mathrm{ref}}\triangleq
\mathbb E_R\{Z_{\mathrm{ref}}-\tau(X)\}^2\).

\begin{theorem}[Finite-class honest validation]
\label{thm:cate-finite-validation}
Let \(M\triangleq|\mathfrak C_j|<\infty\). Suppose the centered RCT and OBS
per-observation scores have range lengths at most \(B_R\) and \(B_O\). Define
\begin{align*}
\varepsilon_{\mathrm{val}}(\delta)&\triangleq
B_R\sqrt{\frac{\log(4M/\delta)}{2n_v}}
+B_O\sqrt{\frac{\log(4M/\delta)}{2N_v}}.
\end{align*}
If \(\widehat c\in\arg\min_c\widehat{\mathcal R}^{\mathrm{ref}}_c\), then,
with probability at least \(1-\delta\),
\begin{align}
\mathcal R_{\mathrm{fit}}^j(\widehat c)
&\leq\min_{c\in\mathfrak C_j}\mathcal R_{\mathrm{fit}}^j(c)
+2\varepsilon_{\mathrm{val}}(\delta).
\label{eq:cate-selector-inequality}
\end{align}
The minimum may be replaced by the larger minimum over
\(\mathfrak C_{j,R}\). This is not realized-sample universal dominance.
\end{theorem}

For an evaluation-independent approximate ratio, define
\(d_{r,c}\triangleq\omega_c\mathbb E_O[(r^\dagger-r_0)G_c]\).
If the reference conditional bias is
\(b_{\mathrm{ref}}(X)\triangleq
\mathbb E_R(Z_{\mathrm{ref}}-\tau(X)\mid X)\), define
\(d_{\mathrm{nuis},c}\triangleq
2\mathbb E_R[b_{\mathrm{ref}}\{\tau-\widehat t_c^j\}]\).
Write the corresponding uniform validation drifts as
\begin{align}
\Delta_r&\triangleq\sup_{c\in\mathfrak C_j}|d_{r,c}|,
&
\Delta_{\mathrm{nuis}}&\triangleq
\sup_{c\in\mathfrak C_j}|d_{\mathrm{nuis},c}|.
\label{eq:validation-drift-bounds}
\end{align}
Theorem~\ref{thm:cate-finite-validation} then adds
\(2\Delta_r+2\Delta_{\mathrm{nuis}}\), provided its
concentration envelope holds for the actual weighted scores.

\subsection{Linear sieves and finite-sample risk}
\label{sec:cate-sieve}

Let \(b(x)\in\mathbb R^p\) be fixed or honestly trained and define
\(t_\beta(x)\triangleq b(x)^\top\beta\). Let \(P\succeq0\) and \(\rho\geq0\).
For DRF, define
\begin{align*}
G_R&\triangleq\mathbb P_R^{\mathrm{fit}}(bb^\top),
&G_{O,r}&\triangleq\mathbb P_O^{\mathrm{fit}}(r_0bb^\top),\\
h_\lambda&\triangleq\mathbb P_R^{\mathrm{fit}}(bZ_\lambda),
&d_g&\triangleq\mathbb P_O^{\mathrm{fit}}(r_0bg)
-\mathbb P_R^{\mathrm{fit}}(bg).
\end{align*}
Whenever the matrix is invertible,
\begin{align}
\widehat\beta_{\lambda,\omega}^{\mathrm{DRF}}
&=\{(1-\omega)G_R+\omega G_{O,r}+\rho P\}^{-1}
\{h_\lambda+\omega d_g\}.
\label{eq:linear-sieve}
\end{align}
For RF, define
\begin{align*}
s_{m_\lambda}&\triangleq
\frac{(A-e_0)(Y-m_\lambda)}{e_0(1-e_0)},
&G_{R,\kappa}&\triangleq\mathbb P_R^{\mathrm{fit}}(\kappa bb^\top),\\
h_{m_\lambda}&\triangleq\mathbb P_R^{\mathrm{fit}}(bs_{m_\lambda}),
&d_{g,\kappa}&\triangleq\mathbb P_O^{\mathrm{fit}}(r_0bg)
-\mathbb P_R^{\mathrm{fit}}(\kappa bg).
\end{align*}
Then
\begin{align}
\widehat\beta_{\lambda,\omega}^{\mathrm{RF}}
&=\{(1-\omega)G_{R,\kappa}+\omega G_{O,r}+\rho P\}^{-1}
\{h_{m_\lambda}+\omega d_{g,\kappa}\}.
\label{eq:linear-sieve-rf}
\end{align}
For \(0\leq\omega\leq1\), nonnegative \(r_0\), and \(P\succeq0\), both
Hessians are positive semidefinite. A feasible DRF or RF fit using
\(\widehat r\) requires \(\widehat r\geq0\), a direct Hessian check, or an
explicit projection onto the positive-semidefinite cone.

For either learner, define
\begin{align*}
G&\triangleq\mathbb E_R(bb^\top),
&\gamma&\triangleq\mathbb E_R(b\tau),\\
\beta_p&\triangleq G^{-1}\gamma,
&t_p&\triangleq b^\top\beta_p,
&a_p^2&\triangleq\mathbb E_R(\tau-t_p)^2.
\end{align*}
Assume \(G\succ0\).
Define the unpenalized Hessian and right-hand-side objects by
\begin{align*}
\widehat H_{\mathrm{DRF}}
&\triangleq(1-\omega)G_R+\omega G_{O,r},
&
\widehat h_{\mathrm{DRF}}
&\triangleq h_\lambda+\omega d_g,\\
\widehat H_{\mathrm{RF}}
&\triangleq(1-\omega)G_{R,\kappa}+\omega G_{O,r},
&
\widehat h_{\mathrm{RF}}
&\triangleq h_{m_\lambda}+\omega d_{g,\kappa}.
\end{align*}
For either \(j\), the penalized normal equation is
\begin{align*}
(\widehat H_j+\rho P)\widehat\beta_j&=\widehat h_j.
\end{align*}
Define
\begin{align*}
\widehat S_j&\triangleq\widehat h_j-\widehat H_j\beta_p.
\end{align*}

\begin{theorem}[Exact sieve excess-risk identity]
\label{thm:cate-sieve-identity}
If \(\widehat H_j+\rho P\) is invertible, then
\begin{align}
\widehat\beta_j-\beta_p
&=(\widehat H_j+\rho P)^{-1}
\{\widehat S_j-\rho P\beta_p\},
\label{eq:sieve-coefficient-identity}\\
\mathbb E_R(\widehat t_j-\tau)^2
&=a_p^2+(\widehat\beta_j-\beta_p)^\top
G(\widehat\beta_j-\beta_p).
\label{eq:sieve-risk-identity}
\end{align}
No same-sample bias-variance split follows from
\(\mathbb E(\widehat S_j)=0\), because the inverse Hessian and score are
dependent.
\end{theorem}

\begin{theorem}[Fixed-candidate finite-sample bound]
\label{thm:cate-fixed-candidate-bound}
Let fitting sizes be \(n_f,N_f\). Let \(\psi_{R,j}\) and \(\psi_{O,j}\)
denote the raw RCT and OBS vector contributions to \(\widehat S_j\), including
their signs and \(\omega\)-multipliers. Define
\begin{align*}
\xi_{R,j}&\triangleq G^{-1/2}\{\psi_{R,j}-\mathbb E_R\psi_{R,j}\},
&
\xi_{O,j}&\triangleq G^{-1/2}\{\psi_{O,j}-\mathbb E_O\psi_{O,j}\}.
\end{align*}
Suppose every one-dimensional projection of these centered whitened scores is
sub-Gaussian with norm at most \(\sigma_{R,j}\) and \(\sigma_{O,j}\),
respectively. On the Gram
event
\begin{align*}
\mathcal E_H&\triangleq
\{\|G^{-1/2}(\widehat H_j-G)G^{-1/2}\|_{\mathrm{op}}\leq\eta\},
\qquad 0\leq\eta<1,
\end{align*}
the following holds with probability at least
\(1-\delta_H-\delta_S\), where \(\Pr(\mathcal E_H^c)\leq\delta_H\):
\begin{align}
\mathbb E_R(\widehat t_j-\tau)^2
&\leq a_p^2+
\frac{2C_{\mathrm{sg}}^2}{(1-\eta)^2}
\left[\sigma_{R,j}\sqrt{\frac{p+\log(4/\delta_S)}{n_f}}
+\sigma_{O,j}\sqrt{\frac{p+\log(4/\delta_S)}{N_f}}\right]^2
\nonumber\\
&\quad+\frac{2\rho^2}{(1-\eta)^2}
\|G^{-1/2}P\beta_p\|_2^2.
\label{eq:finite-sieve-bound}
\end{align}
Here \(C_{\mathrm{sg}}>0\) is a universal constant from Euclidean-norm
concentration for means of independent sub-Gaussian vectors.
The order \(a_p^2+O_p(p/n_f+p/N_f)\) additionally requires bounded score
envelopes, a Gram event bounded away from singularity with probability tending
to one, and ridge bias no larger than that order.
\end{theorem}

With an honest estimated ratio, the conditional score and Hessian drifts are
\begin{align}
d_{r,j}&\triangleq\omega\mathbb E_O[(\widehat r-r_0)b(g-t_p)],
\label{eq:cate-ratio-score-drift}\\
H_{r,j}&\triangleq\omega\mathbb E_O[(\widehat r-r_0)bb^\top].
\label{eq:cate-ratio-hessian-drift}
\end{align}
Normalization alone controls neither object.

\begin{corollary}[Operational additive-spline rate]
\label{cor:cate-spline-rate}
Suppose \(\tau(x)=\sum_{\ell=1}^d\tau_\ell(x_\ell)\), with fixed \(d\) and
smoothness \(s>0\). Consider a deterministic \(K_n\), or a prespecified finite
grid selected by Theorem~\ref{thm:cate-finite-validation}. Let
\(p_n\asymp K_n\), \(a_{p_n}^2=O(K_n^{-2s})\), and
\(\delta_n\downarrow0\). Assume exact \(r_0\), known \(e_0\), and uniformly
along the sequence: (i) the eigenvalues of \(G_{p_n}\) are bounded above and
away from zero; (ii) the centered whitened scores and weighted designs are
sub-Gaussian with bounded norms, while \(r_0\) and \(\kappa\) are nonnegative
and bounded; (iii)
\(K_n+\log(1/\delta_n)=o\{\min(n_f,N_f)\}\); and (iv)
\begin{align*}
\rho_n^2\|G_{p_n}^{-1/2}P\beta_{p_n}\|_2^2
&=O\left(K_n^{-2s}+\frac{K_n}{n_f}+\frac{K_n}{N_f}\right).
\end{align*}
If a finite grid is used, also require the validation remainder in
Theorem~\ref{thm:cate-finite-validation} to be no larger than this target rate.
Then
\begin{align}
\mathbb E_R(\widehat t_j-\tau)^2
&=O_p\left(K_n^{-2s}+\frac{K_n}{n_f}+\frac{K_n}{N_f}\right).
\label{eq:cate-spline-rate}
\end{align}
With \(m_{\mathrm{eff}}\triangleq(1/n_f+1/N_f)^{-1}\) and
\(K_n\asymp m_{\mathrm{eff}}^{1/(2s+1)}\), this becomes
\(O_p\{m_{\mathrm{eff}}^{-2s/(2s+1)}\}\), provided the same conditions hold.
If \(N_f\geq c n_f\), the RCT size determines the generic exponent.
\end{corollary}

\subsection{Implementation boundaries and open theory}
\label{sec:cate-scope}

A fixed or honestly trained frozen neural representation can replace \(b\),
conditional on that representation. End-to-end neural fitting is nonconvex and
is not covered. Estimated-ratio guarantees require uniform score, Hessian, and
validation-loss control. Cross-fitting can recover data efficiency, but
overlapping outer training sets do not justify a naive sum of fold variances.
Neither target preservation nor candidate inclusion gives universal
finite-sample improvement for estimated \((\lambda,\omega)\).
